% ---------------------------------------------------------------------
\section{单举 $B$ 衰变}\label{sec:heavy-inclusive}

单举衰变率与 $R_{e^+e^-}$、DIS 一样是编时乘积前向矩阵元的吸收部分（如 $\Gamma(\bar B\to X_u\ell\bar\nu)\propto\Imag\langle\bar B|\ii\!\int\!\dd^4x\,T\{J^\dagger(x)J(0)\}|\bar B\rangle$），可按 $1/m_b$ 作算符乘积展开（OPE）；在谱的端点 $p_X^2\sim m_b\Lqcd$，末态是一个喷注，一整串被 $\Lqcd/m_b$ 压低的项与领头项同阶，须求和成形状函数（shape function）。对 $\bar B\to X_s\gamma$，在 $\bar B$ 静止系中写 $p_b=m_bv+k$（剩余动量 $k\sim\Lqcd$），光子 $q=\sqrt2E_\gamma n_B$ 沿 $-z$，$p_X=m_bv+k-q$ 沿 $+z$。由 $p_X^2=m_b^2-2m_bE_\gamma+2m_bk^0-2E_\gamma(k^0+k^3)+k^2$ 与 $2E_\gamma=m_b+\order{\Lqcd}$，
\begin{equation}\label{eq:heavy-pX2}
p_X^2=m_b\big(m_b-2E_\gamma+\omega\big)+\order{\Lqcd^2},\qquad \omega\equiv\frac{n_A\cdot k}{n_A\cdot v}=k^0-k^3 :
\end{equation}
喷注只通过 $\omega$ 一个分量感受 $b$ 夸克的费米运动（脚本 \texttt{s3\_\allowbreak inclusive.py}）。

\begin{theorem}[单举 $B$ 衰变的重夸克展开 \tagA]\label{thm:heavy-inclusive-b}
(a)~（局域 OPE）总半轻与总辐射衰变率——更一般地，对强子末态能量作了足够宽的抹平的单举量，如谱的低阶矩——有按 $\Lqcd/m_b$ 的渐近展开：领头项是自由 $b$ 夸克的衰变率（含其 $\as$ 修正），没有 $\order{\Lqcd/m_b}$ 项，$1/m_b^2$ 项由
\begin{equation}\label{eq:heavy-lambda12}
\lambda_1=\frac{\langle\bar B|\bar h_v(\ii D)^2h_v|\bar B\rangle}{2m_B},\qquad
\lambda_2=\frac{\langle\bar B|\bar h_v\frac g2\sigma_{\mu\nu}G^{\mu\nu}h_v|\bar B\rangle}{6m_B}
\end{equation}
决定（$[\ii D^\mu,\ii D^\nu]=\ii gG^{\mu\nu}$；$m_{B^*}^2-m_B^2=4\lambda_2+\order{\Lqcd^3/m_b}$）。例如以极点质量 $m_b$ 表示，
\begin{equation}\label{eq:heavy-Gamma-bu}
\Gamma(\bar B\to X_u\ell\bar\nu)=\frac{G_F^2|V_{ub}|^2m_b^5}{192\pi^3}\Big[1+\order{\as}+\frac{\lambda_1-9\lambda_2}{2m_b^2}+\order{\frac{\as\Lqcd^2}{m_b^2},\frac{\Lqcd^3}{m_b^3}}\Big].
\end{equation}
(b)~（端点区）当 $m_b-2E_\gamma=\order{\Lqcd}$ 时，在 $\Lqcd/m_b$ 的领头幂次、对 $\as$ 的所有阶，
\begin{equation}\label{eq:heavy-HJS}
\frac{\dd\Gamma(\bar B\to X_s\gamma)}{\dd E_\gamma}=2\Gamma_0\,H(m_b,\mu)\int_{2E_\gamma-m_b}^{\bar\Lambda}\dd\omega\;m_b\,J\big(m_b(m_b-2E_\gamma+\omega),\mu\big)\,S(\omega,\mu),
\end{equation}
其中 $\Gamma_0$ 是 $b\to s\gamma$ 的领头阶宽度，硬函数 $H=1+\order{\as}$，$J(p^2,\mu)=\delta(p^2)+\order{\as}$ 是沿 $n_A$ 的夸克喷注函数，$\bar\Lambda=m_B-m_b$，形状函数为
\begin{equation}\label{eq:heavy-shape}
S(\omega)=\frac{1}{2m_B}\Big\langle\bar B(v)\Big|\,\bar h_v\,\delta\Big(\omega-\frac{n_A\cdot\ii D}{n_A\cdot v}\Big)h_v\Big|\bar B(v)\Big\rangle
\end{equation}
（写成双局域形式时两个 $h_v$ 之间是沿 $n_A$ 的类光 Wilson 线）；$\bar B\to X_u\ell\bar\nu$ 的端点区有用同一形状函数写出的类似因子化。
非微扰地位：\cref{asm:intro-wightman}、\cref{asm:intro-composite}、\cref{asm:intro-af-coeff}（MA0、MA4）；端点（形状函数）区的因子化另需\cref{asm:intro-bridge}（MA1、MA2）。
% errtype: E-S
\uses{thm:dis-ope, thm:scet-decoupling, asm:intro-wightman, asm:intro-composite, asm:intro-af-coeff, asm:intro-bridge}
\end{theorem}

\begin{proofsketch}
(1)~对编时乘积在 HQET 算符上作 OPE；它与 DIS 的 OPE 是同一类短程展开，全阶地位相同（\cref{thm:dis-ope}，[A]）。(2)~领头算符 $\bar b\gamma^\mu b$ 的前向矩阵元由 $b$ 数守恒精确归一；量纲 4 的唯一候选 $\bar h_v\,\ii v\cdot D\,h_v$ 由 HQET 运动方程为零，故没有 $1/m_b$ 修正\cite{ChayGeorgiGrinstein1990,BigiUraltsevVainshtein1992}；量纲 5 的动能与色磁算符给出 $\lambda_{1,2}$。Manohar 与 Wise 对 $b\to u$ 求得 $\Gamma/\Gamma_b=1+5E_b-6K_b-2G_b$，$E_b=K_b+G_b$\cite[式~(5.11)]{ManoharWise1994}；在本书约定下 $K_b=-\frac{\lambda_1}{2m_b^2}$（$b$ 夸克的动能，以 $m_b$ 为单位）、$G_b=-\frac{3\lambda_2}{2m_b^2}$（自旋能量），代入得 $1-K_b+3G_b=1+\frac{\lambda_1-9\lambda_2}{2m_b^2}$，即\cref{eq:heavy-Gamma-bu}；其中 $\frac{\lambda_1}{2m_b^2}$ 就是时间膨胀因子 $\langle m_b/E_b\rangle-1$（$\lambda_1=-\langle\bm p^2\rangle$；脚本 \texttt{s3\_\allowbreak inclusive.py}）。(3)~OPE 只在离开物理割线处可靠，围道在小的强子不变质量处靠近割线，故 (a) 只对总宽度与抹平的量成立（与\cref{ch:ee}类似）。(4)~在端点区，$[\Lqcd/(m_b-2E_\gamma)]^n$ 型的项同阶，求和成形状函数\cite{Neubert1994a,Neubert1994b,BigiEtAl1994}。在 SCET 中 $b\to s\gamma$ 流匹配到 SCET 流，$H$ 是匹配系数的模方；$B$ 中动量 $\sim\Lqcd$ 的组分相对于喷注是超软的，经 BPS 场重定义退耦（\cref{thm:scet-decoupling}），留下沿 $n_A$ 的 Wilson 线，即\cref{eq:heavy-shape}；喷注部分给出 $J$。RG 一致性使 Sudakov 对数可以求和\cite{BauerFlemingLuke2000}；图形层次的同一结构见\cite{KorchemskySterman1994}，喷注函数见\cite{Bauer:2003pi,Becher:2006qw}（此步只给出 E-S 型陈述）。
\end{proofsketch}

\paragraph{地位与缺口}
全阶论证（[A]）。(a) 的地位是微扰 OPE 的地位（\cref{thm:dis-ope}，[A]），(b) 的地位是 SCET 因子化的地位（\cref{rem:scet-assumptions}，模式完备，E-S 型陈述）。用于\emph{固定运动学点}上的强子衰变率需要局域对偶（\cref{conj:ee-duality}(b)，[C]）；它不是定理的前提，而是应用所需的额外输入。以极点质量表示时，$m_b^5$ 带有相对量级 $\Lqcd/m_b$ 的红外重整化子歧义\cite{Bigi:1994em,Beneke:1994sw}；OPE 中没有能吸收它的 $1/m_b$ 项，所以它必须与系数级数中相应的发散相消，改用短距离质量后线性的红外敏感性消失，与没有 $1/m_b$ 项相符\cite{Beneke:1994bc}。升级途径：(a) 随微扰 OPE 升级；(b) 需对端点区证明模式完备。

\begin{example}[端点区的树图检验]\label{ex:heavy-photon}
取 $H=1$、$J=\delta(p^2)$：$m_b\,\delta(m_b(m_b-2E_\gamma+\omega))=\delta(\omega-2E_\gamma+m_b)$，\cref{eq:heavy-HJS}化为 $\frac{\dd\Gamma}{\dd E_\gamma}=2\Gamma_0S(2E_\gamma-m_b)$，光子能量 $E_\gamma=\frac12(m_b+\omega)$ 直接映出 $\omega$ 的分布，积分给出 $\Gamma_0\int S=\Gamma_0$。支撑 $\omega\le\bar\Lambda$\cite{Neubert1994b}把端点从 $m_b/2$ 移到 $m_B/2$。树图层次上 $\int S=1$，$\int\omega S=\langle v\cdot k\rangle=0$（运动方程），$\int\omega^2S=\langle(n_A\cdot k)^2\rangle/(n_A\cdot v)^2=-\frac{\lambda_1}{3}$（由 $\langle k^\mu k^\nu\rangle=\frac{\lambda_1}{3}(g^{\mu\nu}-v^\mu v^\nu)$ 与 $n_A^2=0$），故 $\langle E_\gamma\rangle=m_b/2$，方差 $-\frac{\lambda_1}{12}>0$：费米运动把部分子的 $\delta$ 函数谱展宽成宽度 $\sim\Lqcd$ 的峰。数值上 $m_{B^*}=5.325$\,GeV、$m_B=5.280$\,GeV 给出 $\lambda_2\approx0.12$\,GeV$^2$，\cref{eq:heavy-Gamma-bu}中 $-\frac{9\lambda_2}{2m_b^2}\approx-2.3\%$（$m_b=4.8$\,GeV；脚本 \texttt{s7\_\allowbreak extra.py}）。
\uses{thm:heavy-inclusive-b}
\end{example}
